We might think of cost-constrained BO as the following constrained optimization problem: 
\begin{equation*}
    \begin{split}
    & \min_{\textbf{X} \in 2^\Omega} \min_{\bx \in \textbf{X}} f(\bx) \\
    & \text{subject to } \sum_{\bx \in \textbf{X}} c(\bx) \leq \tau. 
    \end{split}
\end{equation*}
Because our optimization domain is $2^\Omega$, i.e., the power set of $\Omega$, we have a nested minimization problem. 
We assume $f(\bx)$ outputs not only its value, but also its evaluation cost determined by a cost function $c(\bx)$ ---which is also black-box. 
Our goal is to minimize $f(\bx)$ subject to the total evaluation cost not exceeding $\tau$. 
This is a pre-specified upper bound on the total cost, such as compute time, dollars, or energy consumption.\looseness-1

One recent research direction in BO has been the development of nonmyopic BO~\citep{lam2016bayesian, lam2017lookahead, yue2019lookahead, frazier2018, lee2020efficient, jiang2020efficient}, which account for the impact of future evaluations and are thus able to make better decisions. 
These decisions are computed by modeling BO as an MDP and then approximating its optimal policy. 
We aim to leverage this MDP framework to make similarly principled, nonmyopic decisions in the cost-constrained setting. 
We do this by extending the MDP to a CMDP, which takes into account variable evaluation costs. The next decision in this setting is a approximation to the optimal CMDP policy. 

Figure~\ref{fig:bo_example} illustrates the advantage of accounting for the cost of future evaluations with CMDP. 
The objective and cost, which have been carefully chosen, are plotted in black and red respectively, and the bottom row indicates feasible trajectories of length two as the optimization continues. 
The optimization domains are similarly shaded when evaluations become infeasible. 
A greedy approach\footnote{We show EI over EIpu for space's sake; the former performs better than the latter. The low cost on the left part of the domain causes EIpu to evaluate exclusively there.}, seen in the first row, does not account for the remaining budget, and is therefore unable to evaluate the global minimum before its becomes infeasible due to an insufficient budget. 
A nonmyopic policy, seen in the second row, better accounts for the cost of future evaluations; it sees that by the fourth iteration, the right side of the domain becomes infeasible, and decides to evaluate there earlier. As a result, it gets much closer to the global minimum. 

In the next section, we formalize our CMDP framework.
We note that our framework is vaguely related to BO with resources (BOR)~\citep{dolatnia2016bayesian}, who consider a partially observable MDP (POMDP) framework for BO when resource consumption of the objective varies, and when there might be multiple agents that can evaluate the acquisition function in parallel. 

\subsection{BO as a Constrained Markov Decision Process}
Given a deterministic cost function $c(\bx): \Omega \rightarrow \R^+$, a cost budget $\tau$, and a GP prior over the observation set $\D_t$ with mean $\mu_t$ and kernel $k_t$, we model $h$ steps of cost-constrained BO as the following CMDP: $<T, \bbS, \A, P, R, C, \tau>$. 

Here, $T$ is the set of decision epochs $\{0, 1, \ldots, h-1\}$ representing $h$ steps of BO. 
While we might want to use an infinite horizon, e.g., $h = \infty$ to iterate until our cost budget is exhausted, we assume a finite horizon for tractability. 
Our state space is the set of observations reachable from starting state $\D_t$ with $h$ BO steps, and the action space is $\Omega$; actions correspond to sampling a point in $\Omega$. 

The transition probabilities from state $\D_t$ to state $\D_{t+1}$, where $\D_{t+1} = \D_t \cup \{(\bx_{t+1}, y_{t+1})\}$, given an action $\bx_{t+1}$, are defined as:
\begin{align*}
    & P( \D_{t+1} \mid \D_t, \bx_{t+1}) \\
    & \sim \N(\mu^{(t)}(\bx_{t+1} ; \D_t), K^{(t)}(\bx_{t+1}, \bx_{t+1}; \D_t)). 
\end{align*}
In other words, the probability of transitioning from $\D_t$ to $\D_{t+1}$ is the probability of sampling $y_{t+1}$ from the posterior of $\GP(\mu_t, \sigma_t^2)$ at $\bx_{t+1}$.

Given an action and transition to a new state $\D_{t+1}$, our reward function is derived from the the EI criterion~\citep{jones1998ego}. 
Let $y_t^*$ be the minimum observed value in the observed set $\D_t$, i.e., $y_t^* = \min \{ y_0, \dots, y_t\}$. Then our reward is expressed as:
\begin{align*}
    & R(\D_t, \bx_{t+1}, \D_{t+1}) = ( y_t^* - y_{t+1})^+ \\
    & \equiv  \max (y_t^* - y_{t+1},0).
\end{align*}

Our CMDP cost is given by $c(\bx)$. 
We assume that this cost is deterministic and state-independent; it only depends on the action. 
In practice, the cost function may be learned as well. 
We emphasize that we assume a \textit{deterministic} cost function; the algorithms and theory we establish do not extend trivially to stochastic cost functions.
Finally, we assume a positive scalar constraint $\tau$. 
However, we can extend this to a vector-valued constraint. 
For example, in materials design, there might be a finite amount of each constituent component, each with its own budget~\citep{abdolshah2019costaware}.  

The expected total reward and cost of a policy $\pi$ are
\begin{align*}
    V_h^\pi(\D_k) &= \E \bigg[ \sum_{t=k}^{k + h-1} R(\D_t, \pi_t(\D_t), \D_{t+1}) \bigg] \\
    &= \E \bigg[ \sum_{t=k}^{k + h-1} (y^*_t - y_{t+1} )^+ \bigg] \nonumber, \\
    C_h^\pi(\D_k) &= \E \bigg[ \sum_{t=k}^{k + h-1} c(\pi_t(\D_t)) \bigg].
\end{align*}
More intuitively, $V_h^\pi(\D_k)$ is the expected reduction in the objective function using policy $\pi$, and $C_h^\pi(\D_k)$ is the accompanying expected cost. 
We can represent a trajectory though this CMDP as the sequence:
\[
    (\bx_k, y_k), (\bx_{k+1}, y_{k+1}), \dots, (\bx_{k+1}, y_{k+h}).
\]
As our cost is strictly positive, a trajectory $(\bx_k, y_k), (\bx_{k+1}, y_{k+1}), \dots, (\bx_{k+1}, y_{k+\ell})$ is feasible if $\sum_{i=k}^{k+\ell} c(\bx_i) \leq \tau$ for some $\ell \leq h$.



